\documentclass[11pt]{article}

\usepackage[a4paper,margin=1in]{geometry}
\usepackage[T1]{fontenc}
\usepackage[utf8]{inputenc}
\usepackage{lmodern}
\usepackage{microtype}
\usepackage{amsmath,amssymb}
\usepackage{booktabs}
\usepackage{graphicx}
\usepackage{xcolor}
\usepackage{hyperref}
\usepackage{enumitem}

\hypersetup{
  colorlinks=true,
  linkcolor=blue!60!black,
  citecolor=blue!60!black,
  urlcolor=blue!60!black
}

\newcommand{\draftnote}[1]{%
  \begin{quote}\footnotesize\color{gray}#1\end{quote}%
}

\title{Working Title: A Linear Immorality Direction in a Small Language Model}
\author{Author names here}
\date{\today}

\begin{document}
\maketitle

% Intentionally no abstract yet.

\section{Introduction}
\draftnote{
Possible framing:
\begin{itemize}[leftmargin=1.5em]
    \item Mechanistic interpretability aims to identify human-understandable variables inside language models.
    \item Prior work has found linear directions for sentiment and other properties, but morality-like representations are less settled, especially in small open models.
    \item The core question here is not whether Qwen behaves morally in some philosophical sense, but whether a simple, stable, causally relevant linear direction separates ETHICS scenarios labeled acceptable vs unacceptable/wrong.
    \item The paper's main claim can be modest: in Qwen2.5-0.5B-Instruct, a layer-15 residual direction fit from ETHICS short commonsense scenarios is stable, predictive, causally relevant, and only partly explained by sentiment.
\end{itemize}
}

\paragraph{Draft contribution list.}
\draftnote{
Something like:
\begin{itemize}[leftmargin=1.5em]
    \item We identify a simple layerwise mean-difference direction in raw residual activations that separates ETHICS commonsense scenarios.
    \item We show that this direction is strongest around layer 15 and stable under both increasing sample size and bootstrap resampling.
    \item We show causal relevance by ablating the direction and measuring a substantial AUROC drop on moral judgment prompts.
    \item We compare this direction to a negative-sentiment direction learned from SST-2 and show moderate alignment but clear angular separation.
    \item We show that a sentiment-controlled residual morality direction still has causal force, and that random-direction ablations do not produce comparable effects.
\end{itemize}
}

\section{Related Work}
\draftnote{
This section can stay short and tightly scoped. Likely subsections:
\begin{itemize}[leftmargin=1.5em]
    \item Mechanistic interpretability and linear directions.
    \item Sentiment directions and representation engineering.
    \item Morality / ethics benchmarks in language models.
    \item Prior ``moral direction'' work and why the present paper is narrower and cleaner.
\end{itemize}
Use the literature note already in \texttt{notes/morality\_direction\_literature\_note.md}.
}

\subsection{Mechanistic Interpretability and Linear Directions}
\draftnote{
Mention work on sparse autoencoders, representation engineering, sentiment directions, and the broader idea that useful internal variables may be approximately linear in hidden-state space.
}

\subsection{Moral and Ethical Evaluation in Language Models}
\draftnote{
Explain that ETHICS and related datasets are mostly behavioral benchmarks. The novelty here is not a new moral dataset, but using ETHICS as a source of high-consensus scenarios for hidden-state analysis.
}

\subsection{Closest Prior Work}
\draftnote{
Be explicit and honest:
\begin{itemize}[leftmargin=1.5em]
    \item broad claim ``models have internal moral representations'' is not new;
    \item the contribution here is a residual-stream, small-open-model, layerwise, causal, sentiment-controlled study using ETHICS commonsense short scenarios.
\end{itemize}
}

\section{Experimental Setup}
\subsection{Model}
\draftnote{
Current model:
\begin{itemize}[leftmargin=1.5em]
    \item \texttt{Qwen/Qwen2.5-0.5B-Instruct}
    \item 24 layers
    \item layer 15 selected by held-out ETHICS performance from a layer sweep
\end{itemize}
You can motivate the model as a small, cheap, open model suitable for rapid mechanistic experiments, while acknowledging its language/comprehension limits.
}

\subsection{Moral Dataset}
\draftnote{
ETHICS commonsense short scenarios.
\begin{itemize}[leftmargin=1.5em]
    \item Label 1: unacceptable / wrong / immoral in the dataset framing.
    \item Label 0: acceptable / not wrong.
    \item Training counts used in the main run: 2789 immoral, 3872 acceptable.
    \item Held-out eval: 983 immoral, 1126 acceptable.
\end{itemize}
Be careful with terminology. In prose, you can note that the dataset is about consensus moral acceptability judgments rather than deep moral philosophy.
}

\subsection{Sentiment Control Dataset}
\draftnote{
SST-2, with the sign flipped so the learned control direction points toward negative sentiment.
Matched-size train subset:
\begin{itemize}[leftmargin=1.5em]
    \item 2789 negative
    \item 2789 positive
\end{itemize}
This gives a fairer comparison than letting sentiment use a much larger dataset.
}

\subsection{Activation Readout}
\draftnote{
Use the final token of a fixed prompt template, e.g.
\begin{itemize}[leftmargin=1.5em]
    \item ETHICS representation prompt: \texttt{Scenario: \{text\}\\nMoral evaluation:}
    \item sentiment representation prompt: \texttt{Text: \{text\}\\nSentiment:}
\end{itemize}
The rationale is that the final prompt token is a natural readout position where the preceding scenario has been integrated.
}

\section{Methods}
\subsection{Direction Estimation}
\draftnote{
The simplest possible estimator:
\[
d = \mu_{\text{immoral}} - \mu_{\text{acceptable}}
\]
with $d$ normalized to unit length.
Stress the simplicity: no SAE, no nonlinear probe, no multi-token aggregation beyond the final pooled state.
}

\subsection{Layerwise Sweep}
\draftnote{
Describe the layer sweep over raw residual activations. The main result:
\begin{itemize}[leftmargin=1.5em]
    \item best ETHICS eval AUROC around layer 15
    \item nearby mid/late layers also strong
\end{itemize}
Possible figure: layerwise AUROC plot.
}

\subsection{Held-Out Evaluation}
\draftnote{
The basic question is whether a direction fit on the train split separates acceptable vs immoral scenarios on held-out ETHICS examples.
Likely metrics:
\begin{itemize}[leftmargin=1.5em]
    \item AUROC as the primary metric
    \item accuracy as secondary, with a note that it is threshold-sensitive
\end{itemize}
}

\subsection{Causal Intervention}
\draftnote{
Define the intervention clearly.
If $u$ is the learned unit direction and $h$ the hidden state at the prompt boundary token, the removal intervention is:
\[
h' = h - (h \cdot u)u
\]
The judgment prompt scores \texttt{acceptable} vs \texttt{unacceptable}. Explain why AUROC is the right primary outcome: it measures separation, not threshold calibration.
}

\subsection{Sentiment Control}
\draftnote{
Three comparisons belong here:
\begin{itemize}[leftmargin=1.5em]
    \item full immorality direction
    \item pure negative-sentiment direction
    \item immorality residual after orthogonalizing against sentiment
\end{itemize}
Write down the orthogonalization:
\[
u_{\perp} = \frac{u - (u \cdot s)s}{\|u - (u \cdot s)s\|}
\]
where $u$ is immorality and $s$ is the sentiment direction.
}

\subsection{Random-Direction Null}
\draftnote{
Important control: matched-norm random-direction ablations at the same layer and token position.
This addresses the objection that ``any disturbance'' might hurt performance.
}

\subsection{Bootstrap Direction Uncertainty}
\draftnote{
This section can be short but useful. Explain the stratified bootstrap over training activations:
\begin{itemize}[leftmargin=1.5em]
    \item resample each class with replacement at original class size
    \item refit the mean-difference direction
    \item compare the bootstrap direction to the full-data direction via cosine and angle
\end{itemize}
This yields interpretable angular stability estimates for both immorality and negative sentiment, plus the uncertainty of their alignment.
}

\section{Results}
\subsection{A Stable Linear Immorality Direction Exists}
\draftnote{
Core numbers to include:
\begin{itemize}[leftmargin=1.5em]
    \item full ETHICS sweep best layer: 15
    \item eval AUROC around 0.829 for the direction-only classifier
    \item sample-size stabilization already visible as train size increases
\end{itemize}
Good figure candidates:
\begin{itemize}[leftmargin=1.5em]
    \item layerwise AUROC
    \item sample-size stability curve
\end{itemize}
}

\subsection{Removing the Direction Degrades Moral Separation}
\draftnote{
Use the full held-out causal numbers.
Main AUROC result:
\begin{itemize}[leftmargin=1.5em]
    \item baseline: 0.7758
    \item remove immorality direction: 0.6357
\end{itemize}
State plainly that this is evidence of causal relevance, not merely correlation.
}

\subsection{Sentiment Explains Part, But Not All, of the Signal}
\draftnote{
Key numbers:
\begin{itemize}[leftmargin=1.5em]
    \item remove pure sentiment: AUROC 0.6607
    \item remove sentiment-orthogonalized immorality residual: AUROC 0.7407
    \item cosine between full-data immorality and negative-sentiment directions: 0.5219
    \item angle between them: about 58.5 degrees
\end{itemize}
The wording here matters. They are substantially aligned relative to random high-dimensional directions, but robustly distinct.
}

\subsection{Random Directions Do Not Reproduce the Effect}
\draftnote{
Main null result:
\begin{itemize}[leftmargin=1.5em]
    \item morality removal AUROC drop: about 0.140
    \item 32 random-direction removals: mean drop near 0, max drop about 0.007
\end{itemize}
This is one of the strongest controls in the paper.
}

\subsection{Bootstrap Stability}
\draftnote{
Key numbers:
\begin{itemize}[leftmargin=1.5em]
    \item immorality direction stays within roughly 6 degrees of the full-data estimate under bootstrap resampling
    \item negative-sentiment direction also within roughly 6 degrees
    \item the alignment between them is itself stable
\end{itemize}
Use the note in \texttt{notes/bootstrap\_uncertainty.md}.
}

\subsection{Visualization and Error Analysis}
\draftnote{
Possible figure set:
\begin{itemize}[leftmargin=1.5em]
    \item immorality vs residual PC1 scatter
    \item same plot with negative-sentiment comparison arrow
\end{itemize}
Possible analysis paragraph:
\begin{itemize}[leftmargin=1.5em]
    \item highest-projection immoral examples are extreme violence / abuse / fraud
    \item high-projection acceptable examples often contain violent or taboo lexical content
    \item low-projection immoral examples are often milder, oddly phrased, or linguistically noisy
\end{itemize}
This suggests the direction captures real moral signal, but is strongly shaped by vivid harm-related language.
}

\subsection{Latent-Output Mismatch}
\draftnote{
Include the result that the latent immorality projection aligns with ETHICS labels slightly better than Qwen's explicit one-word judgment.
This is a nice payoff because it suggests that the latent direction can be cleaner than the model's expressed answer in a weak model.
}

\section{Discussion}
\subsection{What Does the Direction Capture?}
\draftnote{
Good discussion point: the direction is neither a full moral theory nor just sentiment.
It seems to combine:
\begin{itemize}[leftmargin=1.5em]
    \item harm / violence
    \item taboo / crime semantics
    \item negative affect
    \item some additional moral signal beyond pure sentiment
\end{itemize}
}

\subsection{Limitations}
\draftnote{
Likely limitations:
\begin{itemize}[leftmargin=1.5em]
    \item small model with clear comprehension limitations
    \item one dataset and one prompt format
    \item linear direction is only an approximation
    \item ETHICS captures consensus labels, not universal morality
    \item residual-stream analysis alone does not identify a full circuit
\end{itemize}
}

\subsection{Future Work}
\draftnote{
Natural next steps:
\begin{itemize}[leftmargin=1.5em]
    \item repeat on Gemma
    \item extend to controversial items and latent/output comparisons
    \item connect the direction to SAE features or circuits
    \item compare pretrained vs instruct-tuned models
\end{itemize}
}

\section{Conclusion}
\draftnote{
Keep the conclusion modest and crisp.
Something like:
\begin{itemize}[leftmargin=1.5em]
    \item there is a simple, stable, causally relevant linear direction in Qwen 0.5B that tracks ETHICS moral wrongness / unacceptability judgments
    \item this direction is related to but not reducible to negative sentiment
    \item small models already contain morally relevant latent variables that can be studied mechanistically
\end{itemize}
}

\appendix

\section{Prompt Templates}
\draftnote{
Put all representation and judgment templates here.
}

\section{Additional Tables}
\draftnote{
Useful appendix tables:
\begin{itemize}[leftmargin=1.5em]
    \item layerwise metrics
    \item causal AUROC comparison table
    \item bootstrap summary table
\end{itemize}
}

\section{Example Error Cases}
\draftnote{
Pull examples from \texttt{notes/error\_modes.md} and \texttt{notes/latent\_expressed\_mismatch.md}.
}

% Bibliography can be added later.
% \bibliographystyle{plainnat}
% \bibliography{references}

\end{document}
